\documentclass[10pt,a4paper]{article}
\usepackage[T1]{fontenc}
\usepackage{lmodern}
\usepackage[margin=0.8in,top=0.7in,bottom=0.75in]{geometry}
\usepackage{amsmath,amssymb}
\usepackage{booktabs,array,graphicx,float,xcolor,microtype,enumitem,titlesec,listings,caption,tabularx}
\usepackage[hidelinks]{hyperref}
\usepackage{fancyhdr}
\graphicspath{{../figures/}{figures/}}

\definecolor{accent}{HTML}{1F4E79}
\definecolor{boxbg}{HTML}{EEF3F8}
\titleformat{\section}{\large\bfseries\color{accent}}{\thesection}{0.6em}{}
\titleformat{\subsection}{\normalsize\bfseries}{\thesubsection}{0.6em}{}
\titlespacing*{\section}{0pt}{9pt}{4pt}
\titlespacing*{\subsection}{0pt}{6pt}{2pt}
\setlength{\parskip}{3pt}\setlength{\parindent}{0pt}
\setlist{nosep,leftmargin=1.4em}
\captionsetup{font=small,labelfont=bf,skip=3pt}
\pagestyle{fancy}\fancyhf{}\renewcommand{\headrulewidth}{0pt}
\fancyfoot[C]{\small\thepage}
\setlength{\textfloatsep}{6pt}\setlength{\floatsep}{5pt}\setlength{\intextsep}{5pt}

\newcommand{\takeaway}[1]{\par\vspace{2pt}\noindent\colorbox{boxbg}{\parbox{\dimexpr\linewidth-2\fboxsep}{\small\textbf{Takeaway.} #1}}\par\vspace{2pt}}
\newcommand{\fig}[3]{\begin{figure}[H]\centering\includegraphics[width=#2\linewidth]{#1}\caption{#3}\end{figure}}

\lstset{basicstyle=\ttfamily\fontsize{7.4}{8.8}\selectfont,breaklines=true,showstringspaces=false,frame=single,rulecolor=\color{gray!50},
 keywordstyle=\color{accent}\bfseries,commentstyle=\color{gray},language=Python,columns=fullflexible,xleftmargin=0pt,framexleftmargin=2pt}

\IfFileExists{../pidlab/pid.py}{\newcommand{\pidpath}{../pidlab/pid.py}}{\newcommand{\pidpath}{pidlab/pid.py}}
\begin{document}

\begin{center}
{\LARGE\bfseries PID Control Explained}\\[3pt]
{\large Theory, simulation, real-world problems and tuning}\\[4pt]
Vichruth \quad|\quad Modern Monk Labs Technical Assignment \quad|\quad Code: \texttt{github.com/vichruth/ModernMonks-PID-Control}
\end{center}
\vspace{-4pt}\hrule\vspace{5pt}

\noindent\textbf{What this document does.} It explains PID control from scratch and backs every claim with a simulation written in plain NumPy and Matplotlib (no control library). The test machine is a 1\,kg block on a spring and damper that must move to 1\,m and stay there. Every figure and number comes from the code in the repository; \texttt{python run\_all.py} regenerates all of them.

\noindent\textbf{How I used AI.} I used AI tools (Gemini and Claude) as pair-programming help for code scaffolding, plotting and proofreading, and for an independent review: a second AI re-ran key results with its own simulation and caught several wrong numbers, which are corrected here. Results that matter (e.g.\ the steady-state error formula, damping ratio, Ziegler--Nichols $K_u$ and $T_u$) were also checked against closed-form theory.

%==================================================================
\section{The basic idea}
Think of adjusting a shower. The water is too cold, so you open the hot tap a lot. As it warms up you open it less. If it gets too hot you close it again. You never ``know'' the right tap position; you keep measuring the gap between what you want and what you have, and you act on that gap. This is \emph{feedback control}.

In engineering words: the \textbf{setpoint} $r$ is the target, the \textbf{plant} is the machine, a \textbf{sensor} measures its output $y$, and the \textbf{error} is
\begin{equation}e(t)=r(t)-y(t).\end{equation}
The \textbf{controller} turns the error into a command $u$ (here a force in newtons) for the \textbf{actuator}. Without feedback (``open loop'') the machine just runs a fixed command and cannot react to anything unexpected.

\fig{fig01_loop_diagram.png}{0.78}{The control loop: target $\to$ error $\to$ controller $\to$ machine $\to$ sensor $\to$ back to the start. The shower analogy is mapped onto each box.}

\emph{Worked number.} Stepping into the shower at $y=20^\circ$C with a target of $40^\circ$C gives $e=+20^\circ$C, so a big correction. At $39.5^\circ$C the error is $+0.5^\circ$C and almost no correction is needed.
\takeaway{A controller keeps shrinking the gap between target and reality by acting on the measured error.}

%==================================================================
\section{The PID formula}
A PID controller adds up three reactions to the error:
\begin{equation}
u(t)=\underbrace{K_p\,e(t)}_{\text{P: present error}}+\underbrace{K_i\!\int_0^t e(\tau)\,d\tau}_{\text{I: accumulated past error}}+\underbrace{K_d\,\frac{de(t)}{dt}}_{\text{D: current rate of change}}
\end{equation}
\begin{center}\small
\begin{tabularx}{\linewidth}{@{}l l l X@{}}
\toprule Symbol & Name & Unit & Meaning\\\midrule
$r,\;y,\;e$ & setpoint, output, error & m & target position, measured position, $r-y$\\
$u$ & control force & N & what the actuator pushes with\\
$K_p$ & proportional gain & N/m & like a spring: push proportional to how far away you are\\
$K_i$ & integral gain & N/(m$\cdot$s) & memory: keeps growing the push while any error remains\\
$K_d$ & derivative gain & N$\cdot$s/m & like a damper: pushes against fast motion, braking the approach\\
\bottomrule\end{tabularx}\end{center}

\textbf{Plain-words meaning.} P reacts to \emph{how big} the error is now. I reacts to \emph{how long} it has been there. D reacts to \emph{how fast} the error is changing, so it brakes when the target is approached quickly. D does not see the future: it reacts to the current slope, which hints at where the error is heading.

\textbf{What a computer actually does.} It samples every $\Delta t$ (1\,ms here) and repeats:
\begin{align}
e_k &= r_k-y_k & I_k &= I_{k-1}+K_i\,e_k\,\Delta t\\
P_k &= K_p\,e_k & D_k &= K_d\,\frac{e_k-e_{k-1}}{\Delta t},\qquad u_k=P_k+I_k+D_k
\end{align}
At $t=0$ the setpoint steps from 0 to 1\,m, so $P(0)=K_p\cdot 1$\,N; the integral starts at 0 and grows at $K_i\cdot1$\,N/s; $D$ appears only once the block moves.

\fig{fig02_pid_terms.png}{0.55}{Top: the error shrinking over time. Middle: the separate P, I and D forces. Bottom: their sum, the force actually applied.}
\takeaway{P acts on the present, I on the past, D on the trend. The sum of the three is the force.}

%==================================================================
\section{The test system (given)}
A block of mass $m=1$\,kg sits on a spring ($k=20$\,N/m) with a damper ($c=10$\,N$\cdot$s/m). With force $u$:
\begin{equation}m\ddot x+c\dot x+kx=u.\end{equation}
The open-loop damping ratio is $\zeta=c/(2\sqrt{mk})=10/(2\sqrt{20})=1.118>1$, so on its own it creeps to its final position without bouncing. Poles: $s=-2.76$ and $-7.24$.

\textbf{Why a constant push is needed.} Once the block rests at $x=1$\,m, velocity and acceleration are zero, so $0+0+k\,x=u$ and
\begin{equation}u_{ss}=k\,x_{ss}=20\ \text{N/m}\times1\ \text{m}=\mathbf{20\ N}.\end{equation}
The spring keeps pulling the block back, so something must keep pushing with 20\,N. With no push the block returns to 0; with 10\,N it settles at 0.5\,m (Figure~3).

\fig{fig03_open_loop_plant.png}{0.74}{Open-loop responses to constant forces of 0, 10 and 20\,N. Only 20\,N holds the block at 1\,m.}

%==================================================================
\section{How to judge a result}
Four numbers describe a step response (all measured on the target $r=1$\,m):
\begin{enumerate}
\item \textbf{Rise time} $t_r$: time to go from 10\% to 90\% of the target. Reported as ``n/a'' if the response never reaches 90\%.
\item \textbf{Overshoot}: $\text{OS}=\dfrac{y_{\text{peak}}-r}{r}\times100\%$, how far past the target it goes (always measured against the setpoint).
\item \textbf{Settling time} $t_s$: the time after which $y$ stays within $\pm2\%$ of the target (0.98 to 1.02\,m). Reported as ``never'' if it does not.
\item \textbf{Steady-state error} $e_{ss}$: the remaining gap after a long run ($T=15$\,s, averaged over the last 5\% of samples).
\end{enumerate}
\fig{fig04_step_metrics.png}{0.74}{All four measures marked on one example response (PID with $K_p=220$, $K_i=180$, $K_d=14$, chosen only to illustrate the definitions): rise 0.163\,s, overshoot 1.4\%, settling 0.247\,s, steady error 0.0004\,m.}
\takeaway{A good controller is fast, barely overshoots, settles quickly and ends exactly on target, without demanding an impossible force.}

%==================================================================
\section{Each part alone}
\subsection{P only: it always stops short}
With $u=K_p(r-x)$ the closed loop is $m\ddot x+c\dot x+(k+K_p)x=K_p r$. At rest this gives $(k+K_p)x_{ss}=K_p r$, so
\begin{equation}x_{ss}=\frac{K_p}{k+K_p}\,r,\qquad e_{ss}=\frac{k}{k+K_p}\,r.\end{equation}
\emph{In words:} P pushes only while there is an error, and the spring needs a push to hold the block. The block stops where the P-push $K_p\,e_{ss}$ equals the spring force $k\,x_{ss}$; that is always short of the target. For $K_p=10$ the error is $20/30=0.667$\,m. A bigger $K_p$ shrinks the gap but never closes it, and it lowers the damping ratio $\zeta=c/(2\sqrt{m(k+K_p)})$, so the block starts to bounce.

\fig{fig05_p_only.png}{0.78}{P-only step responses for $K_p=10,50,100,300,1000$. Higher gain shrinks the offset (dotted lines are the predicted final values) but adds bouncing and demands a large force.}
\fig{fig06_p_theory_vs_sim.png}{0.78}{Theory (line) against simulation (dots). Left: steady error $k/(k+K_p)$. Right: overshoot against gain.}

\begin{center}\small
\begin{tabular}{@{}rcccccc@{}}
\toprule $K_p$ (N/m) & $\zeta$ & theory $e_{ss}$ (m) & simulated $e_{ss}$ (m) & overshoot & rise time & settling\\\midrule
10 & 0.913 & 0.6667 & 0.6667 & 0\% & n/a & never\\
50 & 0.598 & 0.2857 & 0.2857 & 0\% & n/a & never\\
100 & 0.456 & 0.1667 & 0.1667 & 0\% & 0.184\,s & never\\
300 & 0.280 & 0.0625 & 0.0625 & 31.2\% & 0.077\,s & never\\
1000 & 0.157 & 0.0196 & 0.0196 & 57.6\% & 0.036\,s & 1.43\,s\\\bottomrule\end{tabular}\end{center}
Theory and simulation agree to six decimals. Settling is ``never'' whenever the offset itself exceeds 2\%.
\takeaway{P alone cannot remove the error: $e_{ss}=k/(k+K_p)$. Raising $K_p$ trades a smaller offset for more bounce (57.6\% overshoot at $K_p=1000$).}

\subsection{Adding I: memory removes the offset}
The integral keeps growing while any error remains, so the output can only stop changing when $e=0$. At that point P and D are both 0, so \emph{the integral term alone supplies the 20\,N holding force}: $u_I(\infty)=k\,x_{ss}=20$\,N. That is exactly what P alone could not provide.
\fig{fig07_add_i.png}{0.74}{Adding I at $K_p=100$ ($K_i=0$ to 500). Steady error goes to zero for every $K_i>0$ and the integral force settles at 20\,N. Larger $K_i$ also brings more overshoot.}

\subsection{Adding D: damping removes the bounce}
Taking D from the measurement ($u_D=-K_d\dot x$) simply adds damping, giving $\zeta=\dfrac{c+K_d}{2\sqrt{m(k+K_p)}}$. For $K_p=300$, $\zeta=1$ needs $c+K_d=2\sqrt{320}=35.78$, i.e.\ $K_d\ge25.78$\,N$\cdot$s/m. Textbook overshoot $\exp(-\pi\zeta/\sqrt{1-\zeta^2})$ is measured against the \emph{final} value $K_p/(k+K_p)=0.9375$\,m; against the setpoint it is $\big[0.9375\,(1+\text{OS}_{\text{final}})-1\big]\times100\%$.
\fig{fig08_add_d.png}{0.74}{Adding D at $K_p=300$ ($K_d=0$ to 60). Overshoot falls from 31.2\% to 0\% once $\zeta\ge1$.}
\begin{center}\small
\begin{tabular}{@{}rcccc@{}}
\toprule $K_d$ (N$\cdot$s/m) & $\zeta$ & OS vs final (theory) & OS vs setpoint (theory) & OS vs setpoint (simulated)\\\midrule
0 & 0.280 & 40.07\% & 31.32\% & 31.22\%\\
10 & 0.559 & 12.03\% & 5.02\% & 4.85\%\\
30 & 1.118 & 0\% & 0\% & 0\%\\
60 & 1.957 & 0\% & 0\% & 0\%\end{tabular}\end{center}
D alone still leaves the 6.25\,cm offset (it has no memory), which is why the three terms are combined.
\takeaway{I removes the steady error by learning the 20\,N holding force; D adds damping and removes the bounce.}

%==================================================================
\section{All together}
\fig{fig09_p_pi_pd_pid.png}{0.78}{P, PI, PD and PID on one plot (top) with the force each demands (bottom).}
\begin{center}\small
\begin{tabular}{@{}l c c c c c c@{}}
\toprule Controller & $K_p/K_i/K_d$ & rise (s) & overshoot & settling (s) & $e_{ss}$ (m) & peak force (N)\\\midrule
P & 300/0/0 & 0.077 & 31.22\% & never & 0.0625 & 300.0\\
PI (as in Fig.~9) & 100/200/0 & 0.148 & 21.90\% & 0.845 & 0 & 102.2\\
PD & 300/0/30 & 0.305 & 0\% & never & 0.0625 & 300.0\\
\textbf{PID} & \textbf{300/300/30} & 0.206 & \textbf{4.27\%} & \textbf{1.422} & \textbf{0} & 300.3\\
PID, aggressive & 600/300/30 & 0.102 & 1.18\% & 0.157 & 0 & 600.3\\\bottomrule\end{tabular}\end{center}
(For reference, PI at 300/300/0 overshoots 38.7\% and settles in 0.83\,s.)\par Only the full PID has no steady error \emph{and} no large overshoot: P gives speed, I accuracy, D calm.

\textbf{Why $K_p=300$, $K_i=300$, $K_d=30$.} These gains meet all of my targets: overshoot below 5\% (4.27\%), settling below 1.5\,s (1.42\,s), zero final error. $K_d=30$ is the round number just above the critical-damping value 25.8 (so $\zeta=1.12$), and $K_p=300$ keeps the starting force near 300\,N. I also ran a grid search over $K_p\in[200,600]$, $K_i\in[100,500]$, $K_d\in[10,50]$. It shows that faster settings exist: for example 300/200/20 gives 0.46\% overshoot and 0.243\,s settling at the same 300\,N, and 400/200/20 settles in 0.171\,s at 400\,N. So 300/300/30 is a good, simple baseline, not the optimum. $K_p=600$ shows the other side of the trade-off: it settles in 0.157\,s but demands 600\,N at $t=0$, far above a real 30\,N motor, which leads straight into Section~7.
\fig{fig10_final_gains.png}{0.74}{Selected controller: response and force (top); individual P, I and D forces (bottom).}

%==================================================================
\section{Real-world problems}
Ideal simulations assume unlimited force, perfect sensors and instant updates. Four things break that, each shown with a graph and fixed. All use gains 300/300/30.

\subsection{(a) Actuator limit and integral windup}
A real motor saturates (here $|u|\le30$\,N). While it is saturated the block cannot speed up, but the error is still positive, so the integral keeps filling up. When the block reaches the target the stored charge keeps pushing and it overshoots. \textbf{Fix: clamping anti-windup}: stop integrating while the output is saturated \emph{and} the error would push it further into the limit.
\fig{fig11_windup_clamping.png}{0.74}{Windup under a 30\,N limit. Without the fix: 17.6\% overshoot, 2.98\,s settling. With clamping: 0\% overshoot, 1.06\,s.}

\subsection{(b) Sensor noise and the D term}
Real sensors jitter. With Gaussian noise of $\sigma=1$\,mm (seed 42) and $\Delta t=1$\,ms, differencing two samples gives velocity noise of about $\sqrt2\sigma/\Delta t=1.41$\,m/s, which $K_d=30$ turns into
\begin{equation}\text{std}(u_{\text{noise}})\approx K_d\frac{\sqrt2\,\sigma}{\Delta t}=30\times1.414\approx42.4\ \text{N}\end{equation}
(measured: 43.4\,N), larger than the whole 30\,N limit. \textbf{Fix: low-pass filter the derivative} with $\alpha=\Delta t/(\tau+\Delta t)$: $\tau=0.02$\,s gives 1.76\,N and $\tau=0.05$\,s gives 0.90\,N. The price is a small extra lag.
\fig{fig12_noise_filtering.png}{0.74}{Control force with raw D (large chatter) and with the filtered D ($\tau=0.02$\,s).}

\subsection{(c) Setpoint kick}
If D is taken on the error and the target jumps 0 to 1\,m in one sample, $de/dt=1/0.001=1000$\,m/s and $u_D=30\times1000=30{,}000$\,N, so $u(0)=30{,}300$\,N. \textbf{Fix: differentiate the measurement, not the error.} The block cannot jump, so the spike disappears (300\,N, a 100$\times$ reduction).
\fig{fig13_setpoint_kick.png}{0.66}{Setpoint kick: D on error gives a 30,300\,N spike; D on measurement gives 300\,N.}

\subsection{(d) Slow controller updates}
Real controllers update every $\Delta t_{\text{ctrl}}$ and hold the command in between (zero-order hold). That adds a delay of about $\Delta t_{\text{ctrl}}/2$ (a phase lag $\approx\omega\,\Delta t_{\text{ctrl}}/2$), so the controller reacts to old information. I stepped the plant at 1\,ms and varied the controller period. Long runs show the stability limit:
\begin{itemize}
\item below about 60\,ms the response settles (error shrinks over a 60\,s run);
\item around 62\,ms the oscillation grows slowly;
\item from about 64\,ms upward it grows explosively, and at 90\,ms the peak already reaches 160\,m within 3\,s.
\end{itemize}
A short 3\,s test would report a higher, misleading limit (66.5\,ms), because slow growth is invisible over a few seconds. The system is linear (no saturation) here, so this is steadily growing oscillation, not a limit cycle.
\fig{fig14_slow_sampling.png}{0.74}{Controller update periods of 1, 20, 50 and 90\,ms. Small periods behave like a continuous controller; 90\,ms diverges.}

\begin{center}\small
\begin{tabularx}{\linewidth}{@{}l X X X@{}}
\toprule Problem & Without fix & Fix & With fix\\\midrule
(a) saturation & 17.6\% overshoot, 2.98\,s & clamping anti-windup & 0\% overshoot, 1.06\,s\\
(b) noise & 43.4\,N rms force chatter & low-pass D, $\tau=0.02$\,s & 1.76\,N rms\\
(c) setpoint kick & 30,300\,N spike & D on measurement & 300\,N\\
(d) slow updates & diverges at $\ge64$\,ms & update faster than about 60\,ms & stable, smooth\\\bottomrule\end{tabularx}\end{center}

%==================================================================
\section{Choosing the gains}
\subsection{Manual tuning recipe}
\begin{enumerate}
\item Set $K_i=K_d=0$.
\item Raise $K_p$ until the response is fast but starts to overshoot ($K_p=300$: 31.2\% overshoot).
\item Add $K_d$ until the bounce is damped ($K_d=30$: $\zeta=1.12$, overshoot 0\%).
\item Add $K_i$ until the steady error is gone ($K_i=300$: error 0.0625\,m $\to$ 0).
\item If overshoot returns (4.27\% here), lower $K_i$ or raise $K_d$.
\item Re-check with the force limit and with noise (clamping, low-pass on D).
\end{enumerate}
\begin{center}\small
\begin{tabularx}{\linewidth}{@{}l c c c c c X@{}}
\toprule Raise & rise time & overshoot & settling & steady error & noise & example from this report\\\midrule
$K_p$ & falls & rises & varies & falls & little & $K_p$ 100$\to$1000: $e_{ss}$ 0.167$\to$0.020\,m, OS 0$\to$57.6\%\\
$K_i$ & small fall & rises & improves, then worsens & removed & none & PD$\to$PID: $e_{ss}$ 0.0625\,m $\to$ 0\\
$K_d$ & rises & falls & improves up to $\zeta=1$ & none & much worse & $K_d$ 0$\to$30: OS 31.2\%$\to$0\%\\\bottomrule\end{tabularx}\end{center}

\subsection{Ziegler--Nichols on $G(s)=1/(s+1)^3$}
Method: use P only and raise $K_p$ until the output oscillates steadily; call that gain $K_u$ and the oscillation period $T_u$.

\emph{Theory check.} Three lags have phase $-3\arctan\omega=-180^\circ$ at $\omega=\sqrt3$ rad/s, where $|G|=(1/2)^3=1/8$. So $K_u=8$ and $T_u=2\pi/\sqrt3=3.628$\,s.

\emph{Simulation.} Sweeping $K_p$ from 5 to 10 in steps of 0.1: below 8 the oscillation decays, at 8.0 it is sustained, above 8 it grows. Peaks are 3.625\,s apart, so $K_u=8.00$ and $T_u=3.625$\,s (theory 3.628\,s).
\fig{fig15_zn_sustained_oscillation.png}{0.74}{P-only control around $K_u$: decaying, sustained (period $T_u=3.625$\,s) and growing oscillation.}

The classic PID rule is $K_p=0.6K_u$, $T_i=T_u/2$, $T_d=T_u/8$, converted to parallel gains with $K_i=K_p/T_i$ and $K_d=K_pT_d$. It overshoots a lot, so I compared it with Tyreus--Luyben and searched for a better set (target: overshoot below 10\% and settling below 10\,s).
\fig{fig16_zn_vs_improved.png}{0.74}{Classic Ziegler--Nichols, Tyreus--Luyben and my improved gains on $1/(s+1)^3$.}
\begin{center}\small
\begin{tabular}{@{}l c c c c@{}}
\toprule Method & $K_p/K_i/K_d$ & overshoot & settling & rise\\\midrule
Classic ZN PID & 4.80/2.65/2.18 & 52.4\% & 9.90\,s & 0.95\,s\\
Classic ZN PI & 3.60/1.19/0 & 56.1\% & 30.7\,s & 1.05\,s\\
Tyreus--Luyben$^*$ & 3.64/0.46/2.09 & 0.8\% & 16.9\,s & 1.51\,s\\
\textbf{Improved} ($\tau=0.02$\,s) & \textbf{3.60/0.80/3.20} & \textbf{3.1\%} & \textbf{8.67\,s} & 1.68\,s\\\bottomrule\end{tabular}\\[2pt]{\footnotesize $^*$with D filter $\tau=0.02$\,s (0.96\% overshoot unfiltered). Metrics measured over 35\,s so that slow tails are included.}\end{center}
\takeaway{Ziegler--Nichols gives a starting point, not a finished design: here it needs retuning to cut 52\% overshoot to 3\% with settling under 9\,s.}

%==================================================================
\section{The code}
The controller is one small class in \texttt{pidlab/pid.py} using only NumPy. Options (all switchable, so Section~7 can show before and after): output limits \texttt{u\_min/u\_max}; anti-windup \texttt{"clamp"}, \texttt{"none"} or \texttt{"backcalc"}; D on \texttt{"measurement"} or \texttt{"error"}; and a low-pass filter on D (\texttt{d\_filter\_tau}). Units: $e$ in m, $u$ in N, $K_p$ in N/m, $K_i$ in N/(m$\cdot$s), $K_d$ in N$\cdot$s/m, times in s. Below is its \texttt{update()} method, called once per sample.
\lstinputlisting[firstline=55,lastline=134]{\pidpath}
\subsection*{Run on the test system (30\,N limit, all fixes on)}
\fig{fig17_final_demo.png}{0.62}{Final controller ($K_p=300$, $K_i=300$, $K_d=30$, $\tau=0.02$\,s, $|u|\le30$\,N): position, force and the three terms.}
\begin{center}\small\ttfamily
Gains: Kp=300  Ki=300  Kd=30 | clamp, D on measurement, tau=0.02 s | limit +/-30 N\\
rise 0.454 s | overshoot 0.00\% | settling 0.989 s | steady error 0.0000 m | peak force 30.00 N\\
first 5 forces: [30.0, 30.0, 30.0, 30.0, 30.0]
\end{center}
Compared with Section~6 (4.27\% overshoot, rise 0.206\,s), the 30\,N limit makes the rise slower (0.454\,s) and the clamping removes the overshoot. The force limit, not the gains, is what sets the speed.

%==================================================================
\section{Real robots and myths}
\subsection{Where PID is used}
\begin{enumerate}
\item \textbf{Cruise control.} Error: target speed minus wheel speed. Actuator: throttle. Usually PI: integral action holds speed on hills; D is skipped because wheel-speed signals are noisy.
\item \textbf{Drone attitude.} Error: target angle minus IMU angle (outer loop), target rate minus gyro rate (inner loop). Actuator: motor speed controllers. D braking matters, and it must be filtered because motor vibration leaks into the gyro.
\item \textbf{Robot arm joint.} Error: target angle minus encoder angle. Actuator: motor torque. Holding an arm against gravity needs a steady torque, just like the 20\,N spring force here, so PID is often paired with a gravity feedforward term.
\item \textbf{Line-following robot.} Error: sideways offset of the line from a sensor array. Actuator: difference in left/right wheel speed. Usually PD: P steers, D stops it zig-zagging.
\item \textbf{DC motor speed servo.} Error: target minus measured speed. Actuator: PWM voltage. Often a fast inner current loop inside a slower speed loop; I handles friction.
\end{enumerate}
PID is among the most widely used controllers in industry; real systems pair it with feedforward, cascaded loops and filtering.

\subsection{Five myths, tested}
\begin{enumerate}
\item \textbf{``More gain is always better.''} Raising $K_p$ 100$\to$1000 cuts steady error but overshoot goes 0\%$\to$57.6\% and the starting force reaches 1000\,N (Fig.~5).
\item \textbf{``I fixes everything.''} With a 30\,N limit the integral winds up: 17.6\% overshoot and 2.98\,s settling until clamping is added (Fig.~11).
\item \textbf{``D predicts the future.''} D only reacts to the current slope, and it amplifies noise: 1\,mm of sensor noise makes 43\,N of force chatter (Fig.~12).
\item \textbf{``One set of gains fits all conditions.''} The same 300/300/30 gains give 4.27\% overshoot at 1\,kg, 6.9\% at 2\,kg and 4.0\% at 0.5\,kg (run \texttt{playground.py --mass 2}).
\item \textbf{``Ziegler--Nichols gives the final gains.''} It gave 52.4\% overshoot here; retuning brought it to 3.1\% (Fig.~16).
\end{enumerate}

%==================================================================
\section*{Reproducing the results, and what is missing}
\textbf{Reproduce:} clone the repository, \texttt{pip install -r requirements.txt}, then \texttt{python run\_all.py}, which regenerates all 17 figures and the tables in about 40\,s. \texttt{python playground.py} lets you change gains, mass, noise or force limit and prints the four measures. A test script checks the main numbers against theory.\\
\textbf{Method:} semi-implicit Euler with $\Delta t=1$\,ms (halving $\Delta t$ changes overshoot by 0.01 points); noise is Gaussian with fixed seed 42.\\
\textbf{Limits of this work:} everything is simulation of a linear plant; no hardware was used. The chosen gains are a good baseline, not the optimum (Section~6). The slow-update limit (about 60\,ms) depends on the stated test (long-run growth).

\end{document}